% =====================================================================================================================
% Preamble of the FML book. Packages are restricted to a standard TeX Live installation.
% =====================================================================================================================
\PassOptionsToPackage{obeyspaces,spaces}{url}
\usepackage[T1]{fontenc}
\usepackage[utf8]{inputenc}
\usepackage{lmodern}
\usepackage{mathpazo}
\usepackage[a4paper,inner=2.6cm,outer=2.2cm,top=2.6cm,bottom=2.8cm,headheight=14pt]{geometry}
\usepackage{microtype}
\usepackage{amsmath,amssymb,mathtools,bm}
\usepackage{graphicx}
\usepackage{xcolor}
\usepackage{booktabs,tabularx,longtable,multirow}
\usepackage{enumitem}
\usepackage{caption,subcaption}
\usepackage{listings}
\usepackage{upquote}
\usepackage{xurl}
\usepackage[most]{tcolorbox}
\usepackage{fancyhdr}
\usepackage{makeidx}
\usepackage[numbers,sort&compress]{natbib}
\usepackage[colorlinks=true,linkcolor=blue!55!black,citecolor=green!40!black,urlcolor=blue!65!black,
            bookmarksnumbered=true]{hyperref}
\usepackage[nameinlink,capitalise,noabbrev]{cleveref}

\makeindex
\renewcommand{\indexname}{Index of classes, functions and parameters}
\setlist{itemsep=1pt,topsep=3pt}
\setcounter{tocdepth}{1}
\setcounter{secnumdepth}{3}
\captionsetup{font=small,labelfont=bf}

% Figures of the Einstein-Boltzmann solver report are used in place (no copies in this folder)
\graphicspath{{figures/}{../FML/Cosmology/EinsteinBoltzmannSolver/report/figures/}}

% ---- Headers ---------------------------------------------------------------------------------------------------------
\pagestyle{fancy}
\fancyhf{}
\fancyhead[LE]{\small\thepage\quad\leftmark}
\fancyhead[RO]{\small\rightmark\quad\thepage}
\renewcommand{\headrulewidth}{0.3pt}
\fancypagestyle{plain}{\fancyhf{}\fancyfoot[C]{\small\thepage}\renewcommand{\headrulewidth}{0pt}}

% ---- Colours ---------------------------------------------------------------------------------------------------------
\definecolor{fmlblue}{RGB}{30,70,140}
\definecolor{fmlgreen}{RGB}{30,110,60}
\definecolor{fmlorange}{RGB}{190,100,20}
\definecolor{fmlgrey}{RGB}{245,245,248}
\definecolor{codekw}{RGB}{120,30,120}
\definecolor{codecomment}{RGB}{90,110,90}

% ---- Code listings ---------------------------------------------------------------------------------------------------
\lstdefinestyle{fmlcpp}{
  language=C++, basicstyle=\ttfamily\footnotesize, keywordstyle=\color{codekw}\bfseries,
  commentstyle=\color{codecomment}\itshape, stringstyle=\color{fmlorange}, showstringspaces=false,
  columns=fullflexible, keepspaces=true, breaklines=true, tabsize=4,
  morekeywords={auto,constexpr,nullptr,override,using,namespace,template,typename}}
\lstdefinestyle{fmlini}{basicstyle=\ttfamily\footnotesize, commentstyle=\color{codecomment}\itshape,
  morecomment=[l]{\#}, morecomment=[l]{--}, columns=fullflexible, keepspaces=true, breaklines=true}
\lstset{style=fmlcpp}

\newtcblisting{cppcode}[1][]{listing only, breakable, colback=fmlgrey, colframe=fmlblue!40, boxrule=0.4pt,
  arc=1pt, left=4pt, right=4pt, top=2pt, bottom=2pt, title={#1}, fonttitle=\small\bfseries,
  coltitle=white, colbacktitle=fmlblue!70, listing options={style=fmlcpp}}
\newtcblisting{paramfile}[1][]{listing only, breakable, colback=fmlgrey, colframe=fmlgreen!50, boxrule=0.4pt,
  arc=1pt, left=4pt, right=4pt, top=2pt, bottom=2pt, title={#1}, fonttitle=\small\bfseries,
  coltitle=white, colbacktitle=fmlgreen!70, listing options={style=fmlini}}

% ---- Boxes -----------------------------------------------------------------------------------------------------------
% Module summary at the start of a chapter or section: \begin{modulebox}{Cosmology/Background} ... \end{modulebox}
\newtcolorbox{modulebox}[1]{breakable, colback=fmlblue!4, colframe=fmlblue!60, boxrule=0.5pt, arc=1pt,
  title={Module: \texttt{#1}}, fonttitle=\small\bfseries, fontupper=\small}
% Algorithm / implementation notes
\newtcolorbox{implementation}[1][Implementation]{breakable, colback=fmlgreen!3, colframe=fmlgreen!55,
  boxrule=0.5pt, arc=1pt, title={#1}, fonttitle=\small\bfseries, fontupper=\small}
% Validation summary
\newtcolorbox{validation}[1][Validation]{breakable, colback=fmlorange!4, colframe=fmlorange!70, boxrule=0.5pt,
  arc=1pt, title={#1}, fonttitle=\small\bfseries, fontupper=\small}
% Planned content of a chapter that is not yet written (removed when the chapter is written)
\newtcolorbox{planned}[1][Planned content]{breakable, colback=red!2, colframe=red!45!black, boxrule=0.5pt,
  arc=1pt, title={#1}, fonttitle=\small\bfseries, fontupper=\small}
% Pitfalls / scope limits
\newtcolorbox{caveat}[1][Scope and limitations]{breakable, colback=yellow!5, colframe=yellow!50!black,
  boxrule=0.5pt, arc=1pt, title={#1}, fonttitle=\small\bfseries, fontupper=\small}

% Chapter status line: \chapterstatus{written}{2026-11-14}
\newcommand{\chapterstatus}[2]{\noindent{\footnotesize\color{black!55}\textsf{Status: #1 \;---\; last updated #2}}\par\medskip}

% ---- Code-related markup ---------------------------------------------------------------------------------------------
\newcommand{\fml}{\textsf{FML}}
\urlstyle{tt}
\newcommand{\file}[1]{\nolinkurl{#1}}                                  % a file path in the repository (breaks at / _ .)
\DeclareRobustCommand{\idxtt}[1]{\nolinkurl{#1}}          % used inside index entries
\newcommand{\cls}[1]{\texttt{\detokenize{#1}}\index{#1@\idxtt{#1}}}  % a class/function (indexed)
\newcommand{\code}[1]{\ifmmode\text{\nolinkurl{#1}}\else\nolinkurl{#1}\fi}                                  % inline code (breaks at _ ::)
\newcommand{\param}[1]{\code{#1}\index{#1@\idxtt{#1} (parameter)}}   % parameter
\newcommand{\nsp}[1]{\texttt{\detokenize{#1}}}                          % a namespace
\newcommand{\arxiv}[1]{\href{https://arxiv.org/abs/#1}{arXiv:#1}}
\newcommand{\eprint}[2][]{\href{https://arxiv.org/abs/#2}{arXiv:#2}}
\newcommand{\github}[1]{\href{https://github.com/#1}{\texttt{github.com/#1}}}

% ---- Physics notation ------------------------------------------------------------------------------------------------
\newcommand{\dd}{\mathrm{d}}
\newcommand{\Hp}{\mathcal{H}}
\newcommand{\Om}{\Omega}
\newcommand{\gt}{\tilde g}
\newcommand{\Mpc}{\,\mathrm{Mpc}}
\newcommand{\hMpc}{\,h\,\mathrm{Mpc}^{-1}}
\newcommand{\Mpch}{\,h^{-1}\mathrm{Mpc}}
\newcommand{\muK}{\,\mu\mathrm{K}^2}
\newcommand{\eV}{\,\mathrm{eV}}
\newcommand{\Mpl}{M_{\rm pl}}
\newcommand{\vk}{\bm{k}}
\newcommand{\vx}{\bm{x}}
\newcommand{\vq}{\bm{q}}
\newcommand{\Tr}{\operatorname{Tr}}
\newcommand{\order}[1]{\mathcal{O}(#1)}
\renewcommand{\vec}[1]{\bm{#1}}
\setlength{\emergencystretch}{2.5em}
